\chapter{Creio na vida eterna}
\chaptermark{Creio na vida eterna}

Este capítulo aprofunda a profissão de fé na ressurreição da carne e na vida eterna, desdobrando os parágrafos do Catecismo da Igreja Católica (CIC 1020--1065) sobre morte, juízo particular e final, céu, purgatório, inferno e consumação escatológica. Para adultos em processo de iniciação cristã pelo RICA, essa esperança pascal ilumina o catecumenato como preparação para a bem-aventurança eterna.

\section{A Morte Cristã e o Juízo Particular (\S 1020--1022)}

O cristão une sua morte à de Jesus como passagem à vida eterna; a Igreja, nos últimos ritos, oferece perdão, unção e Viático com ``gentil assurance''.

A morte encerra o tempo de aceitação ou rejeição da graça; o Novo Testamento afirma juízo imediato após a morte segundo obras e fé, com destinos diversos para almas.

Cada um recebe retribuição eterna na alma imortal no momento da morte: juízo particular refere a vida a Cristo, com entrada no céu (purificação ou imediata) ou danação eterna.

Particular judgment ocorre instantaneamente, baseado em DS de concílios e papas; almas purificadas ou condenadas aguardam ressurreição corporal.

Fontes confirmam imortalidade imediata, mas não cobrem plenamente 1022 além do essencial; doutrina magisterial sustenta retribuição inicial.

Para adultos a partir dos 20 anos no RICA, meditar morte e juízo no rito quaresmal desperta exame de consciência, preparando confissão para juízo misericordioso na Vigília Pascal.

\section{O Céu (\S 1023--1029)}

Os que morrem em graça perfeita vivem eternamente com Cristo, vendo-o face a face como Ele é, em beatitude trinitária.

Céu é comunhão de vida e amor com Trindade, Maria, anjos e bem-aventurados: realização suprema das aspirações humanas.

Viver no céu é "estar com Cristo", conservando identidade verdadeira no Reino; fontes enfatizam visão beatífica.

Doutrina sustenta felicidade definitiva, mas fontes disponíveis limitam-se a 1023--1025; tradição completa visão como essência celestial.

Fontes não detalham plenamente 1026--1029; relevância centra em união escatológica com Deus.

Na catequese RICA adulta, céu como meta no tempo pascal une neófitos à liturgia, ansiando visão de Deus na Eucaristia como penhor.

\section{O Purgatório (\S 1030--1032)}

Todos em graça mas imperfeitamente purificados atingem salvação eterna após morte, via purificação para santidade celestial.

Purgatório é purificação final dos eleitos, distinta do inferno; doutrina de Florença e Trento, com fogo de 1Cor 3,15.

Prática de orações pelos mortos desde 2Mac 12,46: Eucaristia, esmolas, indulgências, penitências em sufrágio.

Igreja honra memória dos defuntos, purificando-os para visão beatífica; tradição patrística sustenta.

Fontes cobrem integralmente, confirmando intercessão eclesial como ato de caridade.

Para catecúmenos adultos no RICA, orações pelo purgatório no rito de compaixão quaresmal exercita comunhão dos santos, aplicando méritos pascais aos falecidos.

\section{O Inferno (\S 1033--1037)}

Não amar Deus por pecado grave separa eternamente: "inferno" como autoexclusão definitiva por escolha livre.

Jesus fala de "Gehenna" e fogo inextinguível para impenitentes; anjos reunirão malfeitores para condenação eterna.

Igreja afirma existência e eternidade do inferno: almas em pecado mortal descem imediatamente, sofrendo separação de Deus.

Punição principal: separação eterna de Deus, perda da felicidade criada para Ele.

Fontes detalham 1033--1035; doutrina magisterial (DS) corrobora fogo eterno sem obscurecer misericórdia.

Em RICA para adultos não iniciados, advertência ao inferno no scrutinium fortalece conversão, invocando Espírito para perseverança na graça batismal.

\section{O Juízo Final (\S 1038--1041)}

Ressurreição geral precede Juízo Final: justos e injustos revivem; Cristo separa ovelhas de cabritos para vida ou punição eterna.

No Juízo, verdade de cada relação com Deus revela-se; obras boas ou omissões julgadas, expondo consequências últimas.

Cristo retorna em glória no tempo conhecido só pelo Pai; justiça e amor divinos triunfam sobre injustiças e morte.

Juízo Final ilumina criação e salvação, manifestando providência; fontes enfatizam separação final.

Cobertura até 1040; 1041 inferido como consumação histórica em Cristo.

Para adultos RICA, Juízo Final no catecumenato desperta responsabilidade missionária, julgando obras de misericórdia rumo à Parusia.

\section{Os Novos Céus e a Nova Terra (\S 1042--1060)}

Fim dos tempos: Reino de Deus em plenitude após juízo; justos reinam com Cristo glorificado, universo renovado.

Escritura chama ``novos céus e nova terra'': realização do plano unificador em Cristo, sem lágrimas nem morte.

Deus habitará com homens na Jerusalém celestial; antigas coisas passam, nova criação escatológica.

Fontes limitam-se a 1042--1044; relevância em renovação cósmica, mas não cobrem plenamente 1045--1060 sobre corpo ressuscitado.

Doutrina sustenta transformação total, integrando ressurreição corporal à nova terra.

Na mistagogia RICA adulta, novos céus inspiram ecologia integral, unindo neófitos à criação gemendo por redenção pascal.

\section{A Esperança Cristã Escatológica (\S 1061--1065)}

Não tenho fontes diretas para 1061-1065; doutrina sobre esperança litúrgica e profissão creedal baseia-se em tradição, mas fontes disponíveis focam mistério pascal sem estes parágrafos específicos.

Liturgia professa esperança na vida eterna via Credo; fontes precedentes sustentam via escatologia geral.

Relevância: ``Creio na vida eterna'' como ato de esperança na graça perseverante.

Ausência de citações diretas limita detalhes; recomendo consulta ao CIC completo para profundidade.

Para RICA, esperança escatológica no envio neófitos motiva testemunho, ansiando consumação.

\section{Conclusão}

A vida eterna abarca morte transformada, juízos, purgatório, inferno evitado, céu e nova criação. Adultos RICA, professem esperança pascal em sacramentos rumo à visão beatífica.